% EM structure and teaching revision 20260928
\subsection{磁通保持、磁通子与材料临界条件}
理想导体由零电阻导出的结论是既有磁通不随时间改变；超导体的平衡磁响应还包括迈斯纳效应。二者对“先加磁场后降温”和“先降温后加磁场”的响应不能简单等同。
闭合超导环的量子条件约束的是磁通子。若把几何电感与动能电感之和记为 $L_{\mathrm{tot}}$，外磁通为 $\Phi_{\mathrm{ext}}$，简化模型给出
\[
 \Phi_{\mathrm{ext}}+L_{\mathrm{tot}}I=n\Phi_0,\qquad
 W_n=\frac{(n\Phi_0-\Phi_{\mathrm{ext}})^2}{2L_{\mathrm{tot}}},\quad
 \Phi_0=\frac{h}{2e}.
\]
相邻两条能量抛物线在半整数外磁通处相交，比较 $W_n$ 可确定热力学最低能态。但实际扫场时，能垒会使环停留在亚稳态，未必立即跳到最低的一支。
由某个整数 $n$ 算得电流之后，还须核对临界电流、临界磁场和温度条件；若超出材料允许范围，零电阻与给定电感的模型就不能继续使用。把“允许的量子态”“最低能态”与“实验能到达的状态”分开，才不会把整数取整规则误作所有情形下的实际演化规律。

